% ADD SKLEARN STUFF (SEE JEDISCIKIT.PY)
% INCLUDING ONE-HOT ENCODING


\chapter[ML classifiers: Na\"{i}ve Bayes (2 of 3)]
{\huge ML classifiers: Na\"{i}ve Bayes (2 of 3)}
\label{ch:naiveBayes2}

How can we automate the calculations of the previous chapter in Python? Let's
assume that our data is in a Pandas \texttt{DataFrame} called ``\texttt{t}''
(for ``training'') as it was in the last chapter (p.~\pageref{jediTraining}).
Now imagine writing this function:

\index{predict@\texttt{predict()} (Jedi/Sith)}

\begin{Verbatim}[fontsize=\small,samepage=true,frame=single,framesep=3mm]
def predict(face, height, demeanor, lightsaber):
    ...
\end{Verbatim}

\index{Jedi}
\index{Sith}

This function takes as input a new (unlabeled) data point and returns a
prediction for it: either the word ``\texttt{Jedi}'' or the word
``\texttt{Sith}.''

\index{training data}

We'll do most of the work outside the function, actually. We'll read in the
file and then pre-compute all the probabilities that we'll need for our
calculations. This is effectively the training stage: we're ``training'' our
classifier function on the data.

If you think it through, the calculations from Chapter~\ref{ch:naiveBayes1}
imply the need for the following probabilities:

\index{prior (probability)}
\index{conditional probability}

\begin{compactitem}
\item $P(label)$ for every one of the target labels (these are the ``priors'')
\item $P(value\ |\ label)$ for every value in a column and every target label
(the conditionals)
\end{compactitem}

In our case, this amounts to:

\begin{compactitem}
\item Priors:
    \begin{compactitem}
    \item $P(\textrm{Jedi})$
    \item $P(\textrm{Sith})$
    \end{compactitem}

\item Conditionals:
    \begin{compactitem}
    \item $P(\textrm{clean}\ |\ \textrm{Jedi})$ and
$P(\textrm{clean}\ |\ \textrm{Sith})$
    \item $P(\textrm{beard}\ |\ \textrm{Jedi})$ and $P(\textrm{beard}\ |\ \textrm{Sith})$
    \item $P(\textrm{tall}\ |\ \textrm{Jedi})$ and $P(\textrm{tall}\ |\ \textrm{Sith})$
    \item $P(\textrm{medium}\ |\ \textrm{Jedi})$ and $P(\textrm{medium}\ |\ \textrm{Sith})$
    \item $P(\textrm{short}\ |\ \textrm{Jedi})$ and $P(\textrm{short}\ |\ \textrm{Sith})$
    \item $P(\textrm{menacing}\ |\ \textrm{Jedi})$ and $P(\textrm{menacing}\ |\ \textrm{Sith})$
    \item $P(\textrm{calm}\ |\ \textrm{Jedi})$ and $P(\textrm{calm}\ |\ \textrm{Sith})$
    \item $P(\textrm{red}\ |\ \textrm{Jedi})$ and $P(\textrm{red}\ |\ \textrm{Sith})$
    \item $P(\textrm{green}\ |\ \textrm{Jedi})$ and $P(\textrm{green}\ |\ \textrm{Sith})$
    \item $P(\textrm{blue}\ |\ \textrm{Jedi})$ and $P(\textrm{blue}\ |\ \textrm{Sith})$
    \end{compactitem}

\end{compactitem}

\index{probs@\texttt{probs}}
\index{dictionary}

These are all straightforward to compute with Pandas. I find it convenient to
create a dictionary, called something like ``\texttt{probs},'' to store all the
probabilities. To compute the priors, all we have to do is figure out what
fraction of the rows go with each label:

\begin{Verbatim}[fontsize=\small,samepage=true,frame=single,framesep=3mm]
probs = {}
probs['Jedi'] = len(t[t.type=='Jedi'])/len(t)
probs['Sith'] = len(t[t.type=='Sith'])/len(t)
\end{Verbatim}

For the conditionals like ``$P(\textrm{calm}|\textrm{Sith})$'', we ask:
\textit{considering only the rows that have \texttt{Sith} labels}, what
fraction also have \texttt{calm} labels? That boils down to this:

\footnotesize
\begin{Verbatim}[fontsize=\small,samepage=true,frame=single,framesep=3mm]
probs['calm|Sith'] = (len(t[(t.type=='Sith') & (t.demeanor=='calm')]) /
    len(t[(t.type=='Sith')]))
\end{Verbatim}
\normalsize

The other nine conditionals follow exactly the same pattern. (It may strike you
that we're repeating a lot of very similar code here, and making slight edits
to each one, and that there's a better, more general way. You're right. Stay
tuned for Section~\ref{codingNbGenerally}.)

\index{spycam}
\index{security camera}

Now for the actual \texttt{predict()} function. All we have to do is multiply
everything out per the above equations and see which of the two numerators is
greater. The result is in Figure~\ref{fig:predictFunc}. Stare carefully at it
and convince yourself that it does precisely the same calculations we did in
Section~\ref{jediNumerical} (p.~\pageref{jediNumerical}). Note that we don't
use the \texttt{t} variable inside this function at all. We've already computed
all the \texttt{probs}, and once we've done that, the \texttt{DataFrame} and
the data itself ceases to matter. We run the function just by passing it the
values derived from our spycam:

\begin{figure}
\index{predict@\texttt{predict()} (Jedi/Sith)}
\begin{Verbatim}[fontsize=\small,samepage=true,frame=single,framesep=3mm]
def predict(face, height, demeanor, lightsaber):
    jediness = (probs[face+"|Jedi"] *
                probs[height+"|Jedi"] *
                probs[demeanor+"|Jedi"] *
                probs[lightsaber+"|Jedi"] *
                probs["Jedi"])
    sithness = (probs[face+"|Sith"] *
                probs[height+"|Sith"] *
                probs[demeanor+"|Sith"] *
                probs[lightsaber+"|Sith"] *
                probs["Sith"])
    if jediness > sithness:
        return "Jedi"
    else:
        return "Sith"
\end{Verbatim}
\caption{Our Na\"{i}ve Bayes ML \texttt{predict()} function.}
\label{fig:predictFunc}
\end{figure}

\begin{Verbatim}[fontsize=\small,samepage=true,frame=single,framesep=3mm]
print(predict("beard","tall","calm","blue"))
print(predict("clean","tall","menacing","red"))
\end{Verbatim}
\vspace{-.2in}

\begin{Verbatim}[fontsize=\small,samepage=true,frame=leftline,framesep=5mm,framerule=1mm]
'Jedi'
'Sith'
\end{Verbatim}


\section{Estimating confidence}

\index{evidence}

Our \texttt{predict()} function, as written, just gives a binary answer: either
\texttt{Sith} or \texttt{Jedi}. It may be useful to know \textit{how confident}
it is in its prediction for any given input. All we need to do is compute the
correct ratio. Remember that \texttt{jediness} and \texttt{sithness} are just
the \textit{numerators} of probabilities, and that those probabilities have the
same denominator. Therefore, in order to compute the true $P(\textrm{Jedi}\ |\
e)$, where $e$ is all our evidence (calm, beard, and so forth), we just need to
compute:

\vspace{-.2in}
\begin{align*}
P(\textrm{Jedi}\ |\ e) = \frac{\textrm{jediness}}{\textrm{jediness} + \textrm{sithness}}
\end{align*}
\vspace{-.2in}

and similarly, 

\vspace{-.2in}
\begin{align*}
P(\textrm{Sith}\ |\ e) = \frac{\textrm{sithness}}{\textrm{jediness} + \textrm{sithness}}
\end{align*}
\vspace{-.2in}

The denominator drops out when we divide the quantities anyway.

\medskip

So let's enhance our \texttt{predict()} function to return a list of
\textit{two} values: the prediction, and the confidence level (on a scale of
.5 to 1, meaning ``only 50\% confident in our prediction'' all the way to
``absolutely 100\% confident in our prediction.'' To wit:

\begin{Verbatim}[fontsize=\small,samepage=true,frame=single,framesep=3mm]
def predict(face, height, demeanor, lightsaber):
    ...
    if jediness > sithness:
        return ["Jedi", jediness/(jediness+sithness)]
    else:
        return ["Sith", sithness/(jediness+sithness)]
\end{Verbatim}

\index{list}

(Note that we're returning a list here, of the two values.) Taking it for a
spin:

\begin{Verbatim}[fontsize=\small,samepage=true,frame=single,framesep=3mm]
predict("clean","tall","menacing","red")
predict("beard","tall","menacing","blue")
\end{Verbatim}
\vspace{-.2in}

\begin{Verbatim}[fontsize=\small,samepage=true,frame=leftline,framesep=5mm,framerule=1mm]
['Sith', 0.9760957241941071]
['Jedi', 0.7298292697172792]
\end{Verbatim}

So we're incredibly confident that our clean-shaven, tall, menacing,
red-wielder is a Sith, but not nearly so confident that our bearded, tall,
menacing blue-wielder is a Jedi.

\vspace{.2in}
\section{Coding it more generally}
\label{codingNbGenerally}

The above calculation of the probabilities left something to be desired, since
it was hardcoded and repetitive. We can improve on this by writing our code in
a more general way, rather than explicitly referring to problem-specific names
like \texttt{'Jedi'} and \texttt{'demeanor'} and \texttt{'blue'}. Here's one
way to do all the calculations:


\begin{Verbatim}[fontsize=\small,samepage=true,frame=single]
df = pd.read_csv("...")

probs = {}
features = df.columns[:-1]
target = df.columns[-1]

for label in df[target].drop_duplicates():
    probs[label] = len(df[df[target] == label])/len(df)
    for feature in features:
        for value in df[feature].drop_duplicates():
            probs[value+"|"+label] = (
                len(df[(df[target] == label) & (df[feature] == value)]) / 
                    len(df[df[target] == label]))
\end{Verbatim}

Here's we're assuming (1) the last column of the DataFrame is the target, and
(2) all the values in the columns are unique (\textit{e.g.}, that we won't
have an ``\texttt{out-of-state}'' column with some ``\texttt{beard}'' values and
also have a ``\texttt{married}'' column with ``\texttt{beard}'' values. But those
are pretty much our only assumptions.

The \texttt{predict()} method itself will be generalized as follows: instead
of taking exactly four parameters named ``\texttt{face}'',
``\texttt{demeanor}'', \textit{etc.}, it will take a \textit{list} of
parameter values, so that it can be called like this:

\begin{Verbatim}[fontsize=\small,samepage=true,frame=single]
>>> predict(['beard','tall','menacing','blue'])
\end{Verbatim}

Here's one implementation of the function:

\begin{Verbatim}[fontsize=\small,samepage=true,frame=single]
def predict(values):
    scores = []
    top_prob = 0
    top_answer = ""
    for label in df[target].drop_duplicates():
        prob = probs[label]                     # Start with the prior
        for value in values:
            prob *= probs[value+"|"+label]
        scores.append(prob)
        if prob > top_prob:
            top_prob = prob
            top_answer = label
    return [ top_answer, top_prob / sum(scores) ]
\end{Verbatim}

We use the variable \texttt{scores} to keep a list of the numerators for each
target label: this is only necessary to compute the confidence level in the
result.

\section{Measuring performance}

So how well does this classifier work? This depends on how we measure its
``goodness.'' Commonly, we'll simply ask, ``what percentage of the time is it
right?'' and count errors the same in either direction (calling ``Sith'' for an
actual Jedi, or calling ``Jedi'' for an actual Sith).

As we know, we can't test a classifier's performance on the data it's trained
on and expect to get anything accurate. So we'll need to split our data into a
training set and a test set. This can be done easily in Pandas with:

\begin{Verbatim}[fontsize=\small,samepage=true,frame=single]
complete_data_set = pd.read_csv("warriors.csv")
train = complete_data_set.sample(frac=.7)
test = complete_data_set.drop(train.index)
t = train
\end{Verbatim}

This call to \texttt{.sample()} chooses 70\% of the original, complete
DataFrame's rows and copies them to a DataFrame called ``\texttt{train}''. Then
the next line makes ``\texttt{test}'' a DataFrame of all the others (the
remaining 30\%). For convenience, we then copy \texttt{train} into the
\texttt{t} variable so that none of our previous code has to change.

All our probabilities are thus computed based \textit{only} on the training
data (the random 70\% of the data set). We can then use the \texttt{test} set
to compute our accuracy. All we need to do is iterate through the
\texttt{test} set, passing each row to \texttt{predict()} and seeing whether
we get the correct answer:

\begin{Verbatim}[fontsize=\footnotesize,samepage=true,frame=single]
num_correct = 0
for i in range(len(test)):
    row = test.iloc[i]
    p = predict(row['face'],row['height'],row['demeanor'],row['lightsaber'])
    if p[0] == row['type']:
        num_correct += 1

print("We got " + str(num_correct/len(test)*100) +  
    "% correct on the test data.")
\end{Verbatim}

We get a different answer every time we run our program, since it will choose
a different training/test split in each case.

\begin{Verbatim}[fontsize=\small,samepage=true,frame=single]
We got 82.63% correct on the test data.
\end{Verbatim}

\section{Numerical/continuous variables}

In all the examples above, we had categorical features. Computing something
like ``$P(\textrm{calm}|\textrm{Jedi})$'' is thus just a matter of counting the
\texttt{Jedi} rows which (exactly) match the value \texttt{calm}. But what if
we have a numeric feature, like a person's age or salary?

One quick-and-dirty way to deal with numeric features is to ``bin'' them, kind
of like a histogram. Note that this is actually what we already did with the
``\texttt{height}'' column in the Jedi example: instead of heights like 4'9"
and 6'6", we had values like \texttt{short} and \texttt{tall}.

This has obvious drawbacks. For one, we're throwing away all the detail
present in the measurement and treating different values exactly alike. For
another, it's up to us to somewhat arbitrarily choose how many bins to use and
where the demarcation lines are. This has an \textit{ad hoc} feel to it.

Better is to treat the continuous variable as what it really is: a number.
\textbf{Gaussian Na\"{i}ve Bayes} is probably the most common way to deal with
this in a principled fashion: we make the assumption that the numeric feature
is distributed \textbf{normally} (``normal''==``Gaussian''), and estimate its
mean and standard deviation from the training data. Then, to estimate
$P(value|target\_label)$ for any particular test data point's value, we choose
the estimated value of its \textbf{probability density} and use that.

Suppose we had a numeric \texttt{height} feature (measured in inches) instead
of a categorical one. Our data might look like this:

\begin{Verbatim}[fontsize=\small,samepage=true,frame=none]
        face  height  demeanor lightsaber  type
       clean    62.5  menacing      green  Sith
       beard    68.3      calm      green  Jedi
       clean    78.1  menacing        red  Sith
       beard    63.9      calm       blue  Jedi
       beard    61.9      calm       blue  Jedi
       clean    68.2      calm       blue  Sith
       clean    80.1  menacing        red  Jedi
       clean    33.5      calm      green  Jedi
                    . . .
\end{Verbatim}

Our Gaussian Na\"{i}ve Bayes assumption will be that \textit{for each class, the
\texttt{height} feature is distributed normally}. If we take all the
\texttt{Jedi} rows, and calculate their \texttt{height}s' mean and standard
deviation, we get:

\begin{Verbatim}[fontsize=\small,samepage=true,frame=single]
>>> warriors[warriors.type=='Jedi'].height.mean()
    64.0571428571429
>>> warriors[warriors.type=='Jedi'].height.std()
    9.661931704216865
\end{Verbatim}

Performing the same calculation for the \texttt{Sith} rows gives us:

\begin{Verbatim}[fontsize=\small,samepage=true,frame=single]
>>> warriors[warriors.type=='Sith'].height.mean()
    68.3902173913043
>>> warriors[warriors.type=='Sith'].height.std()
    9.319673312968604
\end{Verbatim}

Looks like Sith Lords are a bit taller than Jedi in general, and slightly less
variable.

\vspace{.1in}
Let's say we're doing a prediction for a data point with a height of 69
inches. We'll need to know the value of the density at $x$=69 for both our
estimated Jedi and estimated Sith distributions. Using the SciPy library, we
can get ``frozen'' density functions for each distribution:

\begin{Verbatim}[fontsize=\small,samepage=true,frame=single]
import scipy.stats
jedi_height = scipy.stats.norm(64.0571428571429, 9.661931704216865)
sith_height = scipy.stats.norm(68.3902173913043, 9.319673312968604)
\end{Verbatim}

and then, to calculate the densities, we just call \texttt{.pdf()} on those
objects:

\begin{Verbatim}[fontsize=\small,samepage=true,frame=single]
>>> jedi_height.pdf(69)
    0.036230663486274517
>>> sith_height.pdf(69)
    0.042717943988373371
\end{Verbatim}

Not surprisingly, the Sith number comes out a fair bit higher, since according
to our estimated density, the average Sith's height is taller than the average
Jedi's.

Keep in mind that the ``.0362'' and ``.0427'' numbers here are a bit hard to
interpret. They are \textit{not} ``the probability that a Jedi/Sith is exactly
69 inches tall'' -- that probability, as with all continuous random variables,
is \textit{zero}. Instead, it's ``the value of the density at $x$=69'' which
isn't really in any units that make sense. It's okay, though, since we are
throwing away out denominator and doing relative comparisons anyway.

By the way, there are also other useful things you can call on a frozen SciPy
distribution, including \texttt{.cdf()} to get the value of the cumulative
density function at any point, \texttt{.ppf()} to get quantile,
\texttt{.rvs()} to generate random variates, ordinary statistical things like
\texttt{.median()}, \texttt{.mean()}, \texttt{.std()}, and other more exotic
statistical things like moments, confidence intervals, differential entropy,
survival functions, \textit{etc.} See the documentation for details.

\subsection{Incorporating into the classifier}

Let's add this feature to our classifier. First, the not-so-generally-coded
version, which had lines like this:

\footnotesize
\begin{Verbatim}[fontsize=\small,samepage=true,frame=single]
probs['tall|Sith'] = (len(t[(t.type=='Sith') & (t.demeanor=='tall')]) /
    len(t[(t.type=='Sith')]))
\end{Verbatim}
\normalsize

Since \texttt{height} is now numeric, we need to replace all those
height-related lines with code like this:

\begin{Verbatim}[fontsize=\small,samepage=true,frame=single]
sith_heights = t[t.type == 'Sith'].height
sith_height = scipy.stats.norm(sith_heights.mean(),sith_heights.std())
jedi_heights = t[t.type == 'Jedi'].height
jedi_height = scipy.stats.norm(jedi_heights.mean(),jedi_heights.std())
\end{Verbatim}

``\texttt{sith\_height}'' is now our frozen distribution function. In 
\texttt{predict()}, then, we simply use it as another multiplier:

\begin{Verbatim}[fontsize=\small,samepage=true,frame=single]
def predict(face, height, demeanor, lightsaber):
    jediness = (probs[face+"|Jedi"] *
                probs[demeanor+"|Jedi"] *
                probs[lightsaber+"|Jedi"] *
                probs["Jedi"])
    sithness = (probs[face+"|Sith"] *
                probs[demeanor+"|Sith"] *
                probs[lightsaber+"|Sith"] *
                probs["Sith"])

    jediness *= jedi_height.pdf(height)
    sithness *= sith_height.pdf(height)

    if jediness > sithness:
        return "Jedi"
    else:
        return "Sith"
\end{Verbatim}

Notice we've removed the ``\texttt{probs[height+"|Jedi"]}'' lines,
since they're no longer applicable: \texttt{height} is numeric now. Instead,
we multiply our \texttt{jediness}/\texttt{sithness} by the value of the
corresponding class's estimated density at that point.

To incorporate this numeric variable into the more general version of the
code, we need to intelligently distinguish between numeric and categorical
columns. This can be done via the \texttt{.dtype} accessor on a DataFrame
column:

\begin{Verbatim}[fontsize=\small,samepage=true,frame=single]
>>> print(t.lightsaber.dtype)
    dtype('O')
>>> print(t.height.dtype)
    dtype('float64')
\end{Verbatim}

(The ``\texttt{O}'' stands for ``object,'' which is Pandas' non-intuitive way of
telling you that the column is something non-numeric, like a string.) You can
directly compare the result to a value like ``\texttt{np.dtype("float64")}'',
but the safer and preferred way is to use NumPy's \texttt{.issubdtype()}
function which more intelligently detects any numeric type (ints, floats,
\textit{etc.}, no matter the precision).

Our generic implementation begins with:

\begin{Verbatim}[fontsize=\small,samepage=true,frame=single]
probs = {}
densities = {}
target = df.columns[-1]
numeric_features = [ c for c in df.columns[:-1] 
    if np.issubdtype(df[c].dtype, np.number) ]
categorical_features = [ c for c in df.columns[:-1]
    if not np.issubdtype(df[c].dtype, np.number) ]
\end{Verbatim}

We're creating a \texttt{densities} dictionary, somewhat similar to our
\texttt{probs} dictionary, to store feature-specific information for our
numeric variables (\texttt{probs} will continue to store information for
categorical variables.) Instead of simple numeric values, \texttt{densities}
will store frozen density functions.

We now populate \texttt{densities} (and \texttt{probs}):

\begin{Verbatim}[fontsize=\small,samepage=true,frame=single]
for label in df[target].drop_duplicates():
    rows_with_this_label = df[df[target] == label]
    probs[label] = len(rows_with_this_label)/len(df)
    for feature in numeric_features:
        density = scipy.stats.norm(rows_with_this_label[feature].mean(),
            rows_with_this_label[feature].std())
        densities[feature+"|"+label] = density
    for feature in categorical_features:
        for value in df[feature].drop_duplicates():
            probs[value+"|"+label] = (
                len(df[(df[target] == label) & (df[feature] == value)]) / 
                    len(df[df[target] == label]))
\end{Verbatim}

Our generic \texttt{predict()} function needs to be adjusted somewhat, since
if we just pass a list of values, we can't tell \textit{which} feature a
numeric value is supposed to go with. There are various ways to solve this,
but we'll choose to pass a key/value pair string here (like
\texttt{"height=69.3"}):

\begin{Verbatim}[fontsize=\footnotesize,samepage=true,frame=single]
def predict(values):
    scores = []
    top_prob = 0
    top_answer = ""
    for label in df[target].drop_duplicates():
        prob = probs[label]    # Start with the prior
        for value in values:
            if "=" in value:
                # Numeric feature.
                numeric_feature, numeric_value = value.split("=")
                numeric_value = float(numeric_value)
                prob *= densities[numeric_feature+"|"+label].pdf(numeric_value)
            else:
                # Categorical feature.
                prob *= probs[value+"|"+label]
        scores.append(prob)
        if prob > top_prob:
            top_prob = prob
            top_answer = label
    return [ top_answer, top_prob / sum(scores) ]
\end{Verbatim}

And we call it like this:
\begin{Verbatim}[fontsize=\small,samepage=true,frame=single]
>>> predict(['beard','height=65','menacing','blue'])
    ['Jedi', 0.85863336769968801]
>>> predict(['beard','height=85','menacing','blue'])
    ['Jedi', 0.72742228657164021]
\end{Verbatim}

How's this thing do?

\begin{Verbatim}[fontsize=\small,samepage=true,frame=single]
We got 82.33% correct on the test data.
\end{Verbatim}

Not too shabby.

